\section{Introduction}

For this project, I have selected the Machine Learning and Neural Network template, Deep Learning
Breast Cancer Detection. The central question this project seeks to address is: Can deep learning improve
the accuracy of breast cancer detection from mammography images? In exploring this question, the
project focuses on the application of transfer learning to medical imaging data, with the aim of improving
detection performance while maintaining clinical relevance.(Litjens et al., 2017; Ragab et al., 2022).

There were two primary reasons for choosing this topic, one academic and one personal. Academically,
the Machine Learning and Neural Networks module was the most engaging module I undertook at level 6
so far. This project provides an opportunity to deepen my understanding of deep learning techniques and
to apply them in a realistic, data driven context. Looking ahead, I intend to pursue either postgraduate
study or professional work involving machine learning, and this project allows me to develop skills that
are directly relevant to that goal. On a personal level, breast cancer is a disease with significant real-world
impact. Having lost my mother to cancer that metastasised following delayed treatment, this project holds
personal importance and reinforces my motivation to explore how machine learning might contribute to
earlier and more reliable diagnoses in the future.

Breast cancer is the most common cancer in the UK, accounting for approximately 15\% of all new cancer
cases and around 30\% of cancers diagnosed in women. Each year, over 56,000 new cases are diagnosed
in the UK, and approximately one in seven women will develop breast cancer during their lifetime
(Cancer Research UK, 2023). Globally, breast cancer remains a major public health concern, with 2.3
million diagnoses and 670,000 deaths reported worldwide in 2022 according to the World Health
Organisation (World Health Organization, 2022). These statistics highlight the scale of the problem and
provide strong justification for continued research into improving screening and diagnostic processes.

Cancer is typically classified into four stages, which describe how far the disease has progressed at the
time of diagnosis (National Cancer Institute, 2023). Early-stage cancers are confined to their original
tissue, while later stages indicate a spread to nearby lymph nodes or further away organs. This spread,
known as metastasing, occurs when cancer cells break away from the primary tumour and travel through
the bloodstream or lymphatic system (American Cancer Society, 2022). Metastasised cancer is
significantly more difficult to treat and is associated with poorer outcomes and a higher mortality rate. As
a result, early detection plays a crucial role in improving survival rates. When breast cancer is diagnosed
at an early stage, treatment is more effective and the likelihood of long-term survival is substantially
higher (Cancer Research UK, 2023).

Breast cancer screening programmes aim to identify cancer before symptoms develop; however, the
current screening process has its limitations. In the UK and many other countries, there is a shortage of
specialist breast radiologists, leading to increased workloads and time pressure (NHS England, 2023).
Accuracy can vary depending on experience, clinician fatigue, and breast tissue density, which can
obscure tumours on standard mammograms (Elmore et al., 2015). These factors contribute to both false
positives, which can lead to unnecessary biopsies and anxiety, and false negatives, which delay diagnosis
and treatment (McKinney et al., 2020).

Machine learning, and deep learning in particular, offers the potential to address some of these challenges.
Convolutional neural networks have demonstrated strong performance in image-based classification tasks
by learning complex feature representations directly from raw pixel data (Esteva et al., 2017; Litjens et
al., 2017). In the context of mammography, deep learning models may be able to identify subtle patterns
that are difficult for the human eye to distinguish, particularly in dense breast tissue, which can be a real
issue for a lot of patients (Shen et al., 2017; McKinney et al., 2020). Rather than replacing clinicians, such
systems are increasingly viewed as decision support tools that can assist radiologists by highlighting
suspicious cases and reducing workload pressure (McKinney et al., 2020).

This project will aim to explore the use of deep learning techniques for breast cancer detection using
mammography images. In particular, it will look at whether CNNs, when trained on existing medical
imaging datasets, have the potential to identify patterns associated with breast cancer (Litjens et al.,
2017). The project aims to examine the feasibility of this approach within the context of current screening
practices.
